\chapter{Classification par la monodromie}

\section{La monodromie comme invariant fondamental}

Nous avons établi au chapitre précédent qu'une fibration de Lefschetz $\pi: X \to \CP^1$ de genre $g$ avec $k$ points critiques est entièrement caractérisée, à difféomorphisme près préservant la fibration, par sa représentation de monodromie
\[
\rho: \pi_1(\CP^1 \setminus \{p_1,\dots,p_k\}, b_0) \longrightarrow \Gamma_g.
\]

Puisque $\CP^1$ est simplement connexe, le groupe fondamental du complémentaire des $k$ points est un groupe libre à $k$ générateurs $\gamma_1,\dots,\gamma_k$, où $\gamma_i$ est un lacet qui tourne une fois autour de $p_i$ (et contourne les autres points critiques). La monodromie est donc déterminée par la donnée des $k$ éléments
\[
t_i = \rho(\gamma_i) \in \Gamma_g,
\]
et la condition de compatibilité globale, provenant du fait qu'un lacet entourant tous les points critiques est contractile, s'écrit
\[
t_1 t_2 \cdots t_k = I \quad \text{dans } \Gamma_g.
\]

De plus, d'après le théorème de monodromie locale, chaque $t_i$ est un twist de Dehn le long du cycle évanouissant $\delta_i \subset \Sigma_g$ associé au point critique.

\begin{definition}
Une \emph{factorisation de l'identité en twists de Dehn} est une suite $(t_1,\dots,t_k)$ d'éléments de $\Gamma_g$, chacun étant un twist de Dehn le long d'une courbe simple non séparante, telle que $t_1 \cdots t_k = I$.
\end{definition}

Le problème de classification des fibrations de Lefschetz se ramène donc à un problème combinatoire : classifier les factorisations de l'identité en twists de Dehn à une relation d'équivalence naturelle près.

\section{Équivalence de Hurwitz}

\subsection{Nécessité d'une relation d'équivalence}

La factorisation $(t_1,\dots,t_k)$ n'est pas canonique : elle dépend du choix des lacets $\gamma_i$ autour des points critiques. Si l'on change le système de lacets, la factorisation est modifiée. Il faut donc identifier les factorisations qui correspondent à la même fibration.

\begin{definition}
Soit $(t_1,\dots,t_k)$ une factorisation de l'identité dans un groupe $G$. Un \emph{mouvement de Hurwitz} est l'une des deux opérations suivantes :
\begin{enumerate}
\item Permutation cyclique : $(t_1,\dots,t_k) \mapsto (t_2,\dots,t_k,t_1)$.
\item Tressage : $(t_i, t_{i+1}) \mapsto (t_i t_{i+1} t_i^{-1}, t_i)$ pour un indice $i$.
\end{enumerate}
Deux factorisations sont dites \emph{Hurwitz-équivalentes} si on peut passer de l'une à l'autre par une suite finie de mouvements de Hurwitz et de leurs inverses.
\end{definition}

L'interprétation géométrique est claire : le tressage correspond au choix d'un lacet différent autour de deux points critiques voisins. Les mouvements de Hurwitz engendrent l'action du groupe de tresses sur les factorisations.

\begin{theoreme}
Deux fibrations de Lefschetz sur $\CP^1$ sont isomorphes (par un difféomorphisme préservant la fibration) si et seulement si leurs factorisations associées sont Hurwitz-équivalentes.
\end{theoreme}

\subsection{Propriétés des mouvements de Hurwitz}

Les mouvements de Hurwitz préservent le produit total :
\[
(t_i t_{i+1} t_i^{-1}) \cdot t_i = t_i t_{i+1}.
\]
Ils sont donc compatibles avec la condition $t_1\cdots t_k = I$.

Dans le cas du genre 1 où $\Gamma_1 \cong SL(2,\Z)$, les mouvements de Hurwitz correspondent à des relations de conjugaison dans le groupe modulaire.

\section{Le cas particulier du genre 1}

\subsection{Identification avec $SL(2,\Z)$}

Rappelons que pour une surface de genre 1, le groupe de mapping class s'identifie à $SL(2,\Z)$. Un twist de Dehn le long d'une courbe simple non contractile correspond à une matrice unipotente :
\[
U = \begin{pmatrix} 1 & 1 \\ 0 & 1 \end{pmatrix}, \quad V = \begin{pmatrix} 1 & 0 \\ -1 & 1 \end{pmatrix},
\]
et plus généralement, tout conjugué de $U$ ou $U^{-1}$ (car l'orientation du twist peut être positive ou négative).

Ainsi, une fibration de Lefschetz de genre 1 sur $\CP^1$ est décrite par une factorisation
\[
A_1 A_2 \cdots A_k = I \quad \text{dans } SL(2,\Z),
\]
où chaque $A_i$ est conjugué à $U$ ou à $U^{-1}$.

\subsection{Interprétation des cycles évanouissants}

Dans le tore $\Sigma_1$, une courbe simple non contractile est déterminée par sa classe d'homologie $(p,q) \in \Z^2$ avec $p$ et $q$ premiers entre eux. Le twist de Dehn correspondant est conjugué à $U^{pq}$ ? En réalité, un twist le long d'une courbe de pente $p/q$ correspond à la matrice
\[
\begin{pmatrix} 1 - pq & p^2 \\ -q^2 & 1 + pq \end{pmatrix}
\]
dans une base appropriée. Cette matrice est bien un conjugué de $U$.

\subsection{Relations classiques dans $SL(2,\Z)$}

Le groupe $SL(2,\Z)$ admet une présentation bien connue :
\[
SL(2,\Z) \cong \langle U, V \mid (UV)^6 = I, \; U V U = V U V \rangle,
\]
où $U$ et $V$ sont les matrices données ci-dessus. La relation $(UV)^6 = I$ est centrale dans notre étude.

\section{Résultats de classification}

\subsection{Travaux de Moishezon}

Moishezon a étudié systématiquement les surfaces elliptiques simplement connexes. Il a montré que pour une surface elliptique relativement minimale sans multiple fiber sur $\CP^1$, la monodromie détermine entièrement la surface.

\begin{theoreme}[Moishezon]
Soit $\pi: S \to \CP^1$ une surface elliptique relativement minimale avec une section. Alors le nombre de fibres singulières est un multiple de 12, et la monodromie est donnée par une factorisation de l'identité en $k = 12n$ twists de Dehn, où $n$ est lié à la signature par $\sigma(S) = -n$.
\end{theoreme}

Le cas $n=1$ correspond à la surface $E(1)$ (rationnelle), $n=2$ à une surface K3, $n\ge 3$ à des surfaces de type général.

\subsection{Classification de Kas et Matsumoto}

Kas (pour le cas elliptique) et Matsumoto (pour le cas général) ont donné une classification complète des fibrations de Lefschetz de genre 1 sur la sphère.

\begin{theoreme}[Kas, Matsumoto]
Toute fibration de Lefschetz de genre 1 sur $\CP^1$ dont la monodromie est une factorisation de l'identité en twists de Dehn positifs (i.e. chaque $A_i$ est conjugué à $U$ et non à $U^{-1}$) est réalisable comme une surface elliptique complexe. La classification se fait par le nombre de fibres et leurs types, qui correspondent aux différentes factorisations modulo équivalence de Hurwitz.
\end{theoreme}

\subsection{Notion d'ensemble $(H+C)$-complet}

Des travaux plus récents (Auroux, Donaldson, Katzarkov, Seidel) ont introduit la notion d'ensemble $(H+C)$-complet pour classifier les factorisations de l'identité en twists de Dehn. Cette approche permet notamment de traiter le cas des fibrations sur le disque (avec bord) qui interviennent dans la construction de variétés symplectiques par chirurgie.

\begin{definition}
Un ensemble de courbes $\{\delta_1,\dots,\delta_k\}$ dans $\Sigma_g$ est dit $(H+C)$-complet si les twists correspondants engendrent le groupe de mapping class et satisfont certaines relations d'intersection.
\end{definition}

Cette notion est encore un domaine actif de recherche.

\subsection{Classification sur le disque}

Pour les fibrations de Lefschetz sur le disque $D^2$ (i.e. avec une seule valeur critique à l'intérieur et une fibre régulière au bord), le problème se réduit à l'étude d'un seul twist de Dehn. Ces fibrations élémentaires sont les briques de base pour construire des fibrations sur $\CP^1$ par recollement.

\section{Invariants déduits de la monodromie}

\subsection{Signature et formule de Meyer}

La signature d'une fibration de Lefschetz de genre 1 peut se calculer directement à partir de la monodromie grce à un résultat de Meyer.

\begin{theoreme}[Meyer]
Soit $\rho: \pi_1(\CP^1 \setminus B) \to SL(2,\Z)$ la monodromie d'une fibration de Lefschetz de genre 1. Alors la signature de la 4-variété totale est donnée par
\[
\sigma = -\frac{1}{3} \sum_{i=1}^k \varepsilon_i,
\]
où $\varepsilon_i = \pm 1$ selon que la monodromie locale autour du $i$-ème point critique est un twist positif (conjugué à $U$) ou négatif (conjugué à $U^{-1}$).
\end{theoreme}

Cette formule est remarquable : la signature ne dépend que du signe des twists, pas de la manière dont ils sont conjugués.

\subsection{Caractéristique d'Euler}

La caractéristique d'Euler s'exprime simplement par
\[
\chi = k,
\]
puisque la fibre régulière a $\chi=0$ et chaque point critique contribue $1$. Cette formule est cohérente avec l'expression $\chi = 12 \chi(\mathcal{O}_S)$ pour les surfaces elliptiques complexes via la relation $k = 12 \chi(\mathcal{O}_S)$.

\subsection{Existence de sections}

La question de savoir si la fibration admet une section (une application $\sigma: \CP^1 \to X$ telle que $\pi \circ \sigma = \mathrm{id}$) a une traduction en termes de monodromie. Une section correspond à un point fixe dans chaque fibre préservé par la monodromie. Dans $SL(2,\Z)$, cela signifie que la monodromie peut être relevée à un groupe agissant sur un espace pointé, ce qui impose des conditions sur les matrices.

\section{Exemples de factorisations}

\subsection{Factorisation triviale}

La factorisation vide ( $k=0$ ) correspond à la fibration triviale $\Sigma_1 \times \CP^1$. La monodromie est triviale, la signature est nulle.

\subsection{Factorisation pour $E(1)$}

Pour la surface $E(1)$, on a $k=12$ twists, tous positifs (donc $\varepsilon_i = 1$ pour tout $i$). La formule de Meyer donne
\[
\sigma = -\frac{1}{3} \times 12 = -4.
\]
Or la signature d'une surface rationnelle comme $E(1)$ est bien $-8$ ? Attention, il y a une subtilité : $E(1)$ est souvent considérée avec une orientation opposée en topologie. Nous reviendrons sur ce calcul détaillé au Chapitre 4.

La factorisation explicite peut s'écrire
\[
(A_1)(A_2)\cdots(A_{12}) = I,
\]
où chaque $A_i$ est un conjugué de $U$. De plus, on sait que cette factorisation est Hurwitz-équivalente à
\[
(UV)^6 = I.
\]

\subsection{Factorisations pour les surfaces de Dolgachev}

Les surfaces de Dolgachev $E(1)_{p,q}$ sont obtenues à partir de $E(1)$ par des twists logarithmiques. Leur monodromie fait intervenir des twists négatifs, ce qui modifie la signature et permet d'obtenir des surfaces homéomorphes à $E(1)$ mais non difféomorphes.

\section{Discussion : limites de l'approche par la monodromie}

La monodromie est un invariant puissant, mais elle ne capture pas toute la géométrie :

\begin{itemize}
\item La structure complexe ( $j$-invariant de la fibre, périodes) est perdue. Deux surfaces elliptiques complexes distinctes peuvent avoir la même monodromie topologique.
\item L'existence de fibres multiples (au sens de Kodaira) n'est pas directement visible dans la monodromie topologique si l'on a perturbé les fibres en singularités de Lefschetz.
\item La monodromie ne distingue pas les différentes orientations possibles (une fibration et son opposée ont la même monodromie dans $SL(2,\Z)$ mais des signatures opposées).
\end{itemize}
Néanmoins, pour les problèmes de classification topologique des 4-variétés, la monodromie s'avère être l'outil le plus efficace.

